% c06 原书 p.59–74（PDF p.066–081）
\subsection{凝聚层}
\label{sec:1.4}

设 $(X,\cO_X)$ 是\emph{环层空间}。$\cO_X$-模层，或简称 $\cO_X$-模\idx{oxmo}{OX-模}，是指一个层 $\cF$，使得对 $X$ 的每个开子集 $U$，$\cF(U)$ 是 $\cO_X(U)$-模；并且对开子集的每个包含 $V\subset U$，限制映射 $\cF(U)\to\cF(V)$ 与模结构及环同态 $\cO_X(U)\to\cO_X(V)$ 相容。也就是说，对任意截面 $a\in\cO_X(U)$ 与 $s,t\in\cF(U)$，有
\[
(s+t)|_V=s|_V+t|_V,\qquad (as)|_V=a|_V\cdot s|_V.
\]

设 $\cF$ 与 $\cG$ 是 $\cO_X$-模。态射 $\phi\colon\cF\to\cG$ 是指层的态射，使得对 $X$ 的每个开子集 $U$，$\phi(U)\colon\cF(U)\to\cG(U)$ 是 $\cO_X(U)$-模同态。从 $\cF$ 到 $\cG$ 的态射集合记为 $\Hom_{\cO_X}(\cF,\cG)$。它是 $\cO_X(X)$-模。

若 $\cO_X$ 是与常预层 $\bZ$ 相伴的层，则 $\cO_X$-模无非就是交换群的层。

定义 $\mathcal{H}om_{\cO_X}(\cF,\cG)$ 为 $\cO_X$-模
\[
U\mapsto\Hom_{\cO_X|_U}(\cF|_U,\cG|_U),
\]
并定义 $\cF\otimes_{\cO_X}\cG$ 为与预层
\[
U\mapsto\cF(U)\otimes_{\cO_X(U)}\cG(U)
\]
相伴的 $\cO_X$-模。

对任意 $P\in X$，有
\[
(\cF\otimes_{\cO_X}\cG)_P=\cF_P\otimes_{\cO_{X,P}}\cG_P.
\]

给定一族 $\cO_X$-模 $\cF_i$（$i\in I$），它们在 $\cO_X$-模范畴中的直和是与预层
\[
U\mapsto\bigoplus_{i\in I}\cF_i(U)
\]
相伴的层。当 $X$ 是诺特拓扑空间时，这个预层是层。（为此要用到：诺特拓扑空间的开子集都拟紧。）$\cF_i$（$i\in I$）在 $\cO_X$-模范畴中的直积是层
\[
U\mapsto\prod_{i\in I}\cF_i(U).
\]

设 $(I,\le)$ 是有向集。对 $\cO_X$-模的有向系统 $(\cF_i,\phi_{ij})_{i\in I}$，其在 $\cO_X$-模范畴中的有向极限是与预层
\[
U\mapsto\dlim_i\cF_i(U)
\]
相伴的层。若 $X$ 是诺特拓扑空间，则这个预层是层。若 $(\cF_i,\phi_{ji})_{i\in I}$ 是 $\cO_X$-模的逆向系统，则其在 $\cO_X$-模范畴中的逆向极限是层
\[
U\mapsto\plim_i\cF_i(U).
\]

$\cO_X$-模称为\emph{自由}的\idx{ziyou}{自由}，若它同构于若干 $\cO_X$ 的拷贝的直和。拷贝的个数称为自由 $\cO_X$-模的\emph{秩}\idx{zhi}{秩}。$\cO_X$-模 $\cF$ 称为\emph{局部自由}的\idx{jubuziyou}{局部自由}，若存在 $X$ 的开覆盖 $\{U_i\}_{i\in I}$，使每个 $\cF|_{U_i}$ 是自由的 $\cO_{U_i}$-模。对局部自由 $\cO_X$-模 $\cF$，定义其秩为 $X$ 上的函数 $P\mapsto\rank_{\cO_{X,P}}\cF_P$。它是局部常值函数。秩为 $1$ 的局部自由 $\cO_X$-模称为\emph{可逆}的\idx{keni}{可逆层}。可逆 $\cO_X$-模的张量积可逆。此外，若 $\cL$ 是可逆 $\cO_X$-模，则 $\mathcal{H}om_{\cO_X}(\cL,\cO_X)$ 也可逆，且有典范同构
\[
\cL\otimes_{\cO_X}\mathcal{H}om_{\cO_X}(\cL,\cO_X)\cong\cO_X.
\]
于是可逆 $\cO_X$-模的同构类关于张量积构成一个群。称此群为 \emph{Picard 群}\idx{picardqun}{Picard 群}，记作 $\operatorname{Pic}(X)$。
\begin{tnote}
可逆 $\cO_X$-模就是秩为一的局部自由层，几何上对应 $X$ 上的线丛；\emph{Picard 群} $\operatorname{Pic}(X)$ 即这些线丛的同构类按张量积构成的群，它与上同调群 $H^1(X,\cO_X^*)$ 同构。「可逆」一词来自上面那个同构：$\cL$ 关于张量积有逆元 $\cL^{-1}=\mathcal{H}om_{\cO_X}(\cL,\cO_X)$。[依据：原书 p.60 对 Picard 群的定义；标准结论见 Hartshorne, \emph{Algebraic Geometry}, GTM 52, §II.6]
\end{tnote}

下面是一个例子：设 $X$ 是拓扑空间，$C$ 是复值连续函数的层。则 $(X,C)$ 是环层空间。对 $X$ 上任意复向量丛 $p\colon E\to X$，由
\[
U\mapsto\{s\colon U\to p^{-1}(U)\mid s \text{ 连续且 } ps=\id\}
\]
定义的截面层是局部自由的 $C$-模层。

$\cO_X$-模 $\cF$ 称为\emph{有限展示}的\idx{youxianzhanshi}{有限展示}，若存在 $X$ 的开覆盖 $\{U_i\}_{i\in I}$，使在每个 $U_i$ 上有形如
\[
\cO_{U_i}^{m_i}\to\cO_{U_i}^{n_i}\to\cF|_{U_i}\to0
\]
的正合序列，其中 $\cO_{U_i}^{m_i}$ 与 $\cO_{U_i}^{n_i}$ 分别是秩为 $m_i$ 与 $n_i$ 的自由 $\cO_{U_i}$-模。$\cO_X$-模 $\cF$ 称为\emph{有限型}的\idx{youxianxing}{有限型}，若存在 $X$ 的开覆盖 $\{U_i\}_{i\in I}$，使在每个 $U_i$ 上有形如
\[
\cO_{U_i}^{n_i}\to\cF|_{U_i}\to0
\]
的正合序列，其中 $n_i$ 是某个自然数。
\begin{tnote}
「有限展示」（finite presentation）与「有限型」（finite type）是模论术语在层论中的平移，术语表未收前者，此处按字面直译；后者按术语表译作「有限型」。[来源：本库 SY07-08 术语表, 工程文件/glossary.tsv]
\end{tnote}

\begin{proposition}\label{it:1.4.1}
设 $(X,\cO_X)$ 是环层空间，$\cF$ 与 $\cG$ 是两个 $\cO_X$-模。
\begin{enumerate}[label=(\roman*),leftmargin=2.8em]
  \item 若 $\cF$ 是有限展示的，则对每个 $P\in X$，有
  \[
  (\mathcal{H}om_{\cO_X}(\cF,\cG))_P\cong\Hom_{\cO_{X,P}}(\cF_P,\cG_P).
  \]
  \item 设 $\cF$ 与 $\cG$ 都是有限展示的，$P\in X$。若存在 $\cO_{X,P}$-模同构 $\phi_P\colon\cF_P\xrightarrow{\sim}\cG_P$，则存在 $P$ 的邻域 $U$ 与 $\cO_U$-模同构 $\phi\colon\cF|_U\xrightarrow{\sim}\cG|_U$，使 $\phi$ 在茎上诱导同构 $\phi_P$。
\end{enumerate}
\end{proposition}

\begin{proof}
(i) 我们这样定义 $(\mathcal{H}om_{\cO_X}(\cF,\cG))_P\to\Hom_{\cO_{X,P}}(\cF_P,\cG_P)$：把截面 $\phi\in\Hom_{\cO_X}(\cF,\cG)(U)$ 的芽映到 $\phi\colon\cF|_U\to\cG|_U$ 在茎上诱导的同态 $\phi_P\colon\cF_P\to\cG_P$。取 $P$ 的邻域 $U$，使有正合序列
\[
\cO_U^{m}\to\cO_U^{n}\to\cF|_U\to0.
\]
则有交换图
\begin{center}
\begin{tikzcd}[column sep=1.6em]
0 \arrow[r] & (\mathcal{H}om_{\cO_U}(\cF|_U,\cG|_U))_P \arrow[r] \arrow[d] &
(\mathcal{H}om_{\cO_U}(\cO_U^{n},\cG|_U))_P \arrow[r] \arrow[d] &
(\mathcal{H}om_{\cO_U}(\cO_U^{m},\cG|_U))_P \arrow[d] \\
0 \arrow[r] & \Hom_{\cO_{U,P}}(\cF_P,\cG_P) \arrow[r] &
\Hom_{\cO_{U,P}}(\cO_{U,P}^{n},\cG_P) \arrow[r] &
\Hom_{\cO_{U,P}}(\cO_{U,P}^{m},\cG_P).
\end{tikzcd}
\end{center}
最后两个竖直箭头是同构。故第一个竖直箭头也是同构。

(ii) $\phi_P$ 是 $\Hom_{\cO_{X,P}}(\cF_P,\cG_P)$ 中的元素。由 (i)，存在 $P$ 的邻域 $U$ 与态射 $\phi\colon\cF|_U\to\cG|_U$，使 $\phi$ 在茎上诱导同态 $\phi_P$。设 $\psi_P\colon\cG_P\to\cF_P$ 是 $\phi_P$ 的逆。必要时缩小 $U$，可找到态射 $\psi\colon\cG|_U\to\cF|_U$，使 $\psi$ 在茎上诱导同态 $\psi_P$。因 $\psi_P\phi_P=\id_{\cF_P}$ 且 $\phi_P\psi_P=\id_{\cG_P}$，必要时再缩小 $U$，便得 $\psi\phi=\id_{\cF|_U}$ 与 $\phi\psi=\id_{\cG|_U}$。故 $\phi$ 是同构。
\end{proof}

设 $(f,f^\sharp)\colon(X,\cO_X)\to(Y,\cO_Y)$ 是环层空间的态射，$\cF$ 是 $\cO_X$-模。则 $f_*\cF$ 是 $f_*\cO_X$-模。经由环层态射 $f^\sharp\colon\cO_Y\to f_*\cO_X$，$f_*\cF$ 成为 $\cO_Y$-模。称它为 $\cF$ 的\emph{正像}\idx{zhengxiang}{正像}。设 $\cG$ 是 $\cO_Y$-模。则 $f^{-1}\cG$ 是 $f^{-1}\cO_Y$-模。态射 $f^\sharp\colon\cO_Y\to f_*\cO_X$ 诱导环层态射 $f^{-1}\cO_Y\to\cO_X$。定义 $\cG$ 的\emph{逆像}\idx{nixiang}{逆像}为 $\cO_X$-模
\[
f^*\cG=\cO_X\otimes_{f^{-1}\cO_Y}f^{-1}\cG.
\]
留给读者构造双射
\[
\Hom_{\cO_Y}(\cG,f_*\cF)\cong\Hom_{\cO_X}(f^*\cG,\cF),
\]
它对 $\cF$ 与 $\cG$ 都是函子性的。这表明 $f^*$ 是 $f_*$ 的左伴随。对每个 $P\in X$，有
\[
(f^*\cG)_P\cong(\cO_X\otimes_{f^{-1}\cO_Y}f^{-1}\cG)_P
\cong\cO_{X,P}\otimes_{(f^{-1}\cO_Y)_P}(f^{-1}\cG)_P
\cong\cO_{X,P}\otimes_{\cO_{Y,f(P)}}\cG_{f(P)},
\]
即
\[
(f^*\cG)_P\cong\cO_{X,P}\otimes_{\cO_{Y,f(P)}}\cG_{f(P)}.
\]
\begin{tnote}
$f^{-1}$ 与 $f^*$ 不可混用：$f^{-1}\cG$ 只是把层沿 $f$ 拉回，仍是 $f^{-1}\cO_Y$-模；要得到 $\cO_X$-模，须再与 $\cO_X$ 作张量积，这就是 $f^*\cG$。第二个同构用了 $f^{-1}$ 与茎的交换性：$(f^{-1}\cG)_P=\cG_{f(P)}$。[依据：原书 p.62 中逆像的定义与茎的计算]
\end{tnote}

设 $A$ 是环，$M$ 是 $A$-模。定义 $\cO_{\Spec A}$-模 $M^\sim$ 如下：对 $\Spec A$ 的每个开子集 $U$，$M^\sim(U)$ 由满足以下两个条件的映射 $s\colon U\to\coprod_{\fp\in\Spec A}M_\fp$ 组成：
\begin{enumerate}[label=(\alph*),leftmargin=2.8em]
  \item 对每个 $\fp\in U$，有 $s(\fp)\in M_\fp$。
  \item 对每个 $\fp\in U$，存在 $\fp$ 的含于 $U$ 的邻域 $U_\fp$、元素 $m\in M$ 与 $f\in A$，使对每个 $\fq\in U_\fp$，有 $f\notin\fq$ 且 $s(\fq)=\frac{m}{f}$ 于 $M_\fq$ 中。
\end{enumerate}
对开子集的每个包含 $V\subset U$，定义 $M^\sim(U)\to M^\sim(V)$ 为映射的限制。

注意有 $A^\sim=\cO_{\Spec A}$。

\begin{proposition}\label{it:1.4.2}
\begin{enumerate}[label=(\roman*),leftmargin=2.8em]
  \item 对每个 $\fp\in\Spec A$，有 $(M^\sim)_\fp\cong M_\fp$。
  \item 对每个 $f\in A$，有 $M^\sim(D(f))\cong M_f$。特别地，取 $f=1$，得 $M^\sim(\Spec A)\cong M$。经由同构 $D(f)\cong\Spec A_f$，$M^\sim|_{D(f)}$ 与 $M_f^\sim$ 等同。
  \item $A$-模的序列
  \[
  M'\to M\to M''
  \]
  正合，当且仅当 $\cO_{\Spec A}$-模的序列
  \[
  M'^\sim\to M^\sim\to M''^\sim
  \]
  正合。
  \item 对任意 $A$-模 $M$ 与 $N$，有
  \[
  \Hom_{\cO_{\Spec A}}(M^\sim,N^\sim)\cong\Hom_A(M,N),\qquad
  M^\sim\otimes_{\cO_{\Spec A}}N^\sim\cong(M\otimes_A N)^\sim.
  \]
  若 $M$ 是有限展示的 $A$-模，则
  \[
  \mathcal{H}om_{\cO_{\Spec A}}(M^\sim,N^\sim)\cong(\Hom_A(M,N))^\sim.
  \]
  \item 对一族 $A$-模 $M_i$（$i\in I$），有
  \[
  \bigoplus_i M_i^\sim\cong\Bigl(\bigoplus_i M_i\Bigr)^\sim.
  \]
  \item 对 $A$-模的有向系统 $(M_i,\phi_{ij})_{i\in I}$，有
  \[
  \dlim_i M_i^\sim\cong\bigl(\dlim_i M_i\bigr)^\sim.
  \]
\end{enumerate}
\end{proposition}

\begin{proof}
我们只证 (iv)，其余留给读者。定义
\[
\alpha\colon\Hom_{\cO_{\Spec A}}(M^\sim,N^\sim)\to\Hom_A(M,N)
\]
如下：对每个态射 $\phi\colon M^\sim\to N^\sim$，定义 $\alpha(\phi)$ 为复合
\[
M\cong M^\sim(\Spec A)\xrightarrow{\phi}N^\sim(\Spec A)\cong N.
\]
定义
\[
\beta\colon\Hom_A(M,N)\to\Hom_{\cO_{\Spec A}}(M^\sim,N^\sim)
\]
如下：给定同态 $\phi\colon M\to N$，对每个 $\fp\in\Spec A$，设 $\phi_\fp\colon M_\fp\to N_\fp$ 是 $\phi$ 诱导的同态。对 $M^\sim(U)$ 的每个截面 $s\colon U\to\coprod_{\fp\in\Spec A}M_\fp$，定义 $N^\sim(U)$ 的截面 $\beta(\phi)(s)\colon U\to\coprod_{\fp\in\Spec A}N_\fp$ 为
\[
(\beta(\phi)(s))(\fp)=\phi_\fp(s(\fp))
\]
对每个 $\fp\in U$ 成立。于是得到态射 $\beta(\phi)\colon M^\sim\to N^\sim$。可以验证 $\alpha$ 与 $\beta$ 互为逆。故
\[
\Hom_{\cO_{\Spec A}}(M^\sim,N^\sim)\cong\Hom_A(M,N).
\]
定义态射
\[
\gamma\colon M^\sim\otimes_{\cO_{\Spec A}}N^\sim\to(M\otimes_A N)^\sim
\]
如下：因 $M^\sim\otimes_{\cO_{\Spec A}}N^\sim$ 是与预层
\[
U\mapsto M^\sim(U)\otimes_{\cO_{\Spec A}(U)}N^\sim(U)
\]
相伴的层，只需定义从这个预层到 $(M\otimes_A N)^\sim$ 的态射。对 $M^\sim(U)$ 中的截面 $s\colon U\to\coprod_{\fp\in\Spec A}M_\fp$ 与 $N^\sim(U)$ 中的截面 $t\colon U\to\coprod_{\fp\in\Spec A}N_\fp$，定义 $(M\otimes_A N)^\sim(U)$ 中的截面
\[
\gamma(s\otimes t)\colon U\to\coprod_{\fp\in\Spec A}(M\otimes_A N)_\fp,
\]
使对每个 $\fp\in U$，$\gamma(s\otimes t)(\fp)$ 是 $s(\fp)\otimes t(\fp)$ 在典范同构
\[
M_\fp\otimes_{A_\fp}N_\fp\cong(M\otimes_A N)_\fp
\]
下的像。可以验证 $\gamma$ 在 $\fp$ 处的茎上诱导的同态正是这个典范同构。故 $\gamma$ 是同构。

对每个 $\fp\in\Spec A$，有典范同态
\[
(\Hom_A(M,N))_\fp\to\Hom_{A_\fp}(M_\fp,N_\fp).
\]
设 $M$ 是有限展示的。取正合序列
\[
A^m\to A^n\to M\to0.
\]
则有交换图
\begin{center}
\begin{tikzcd}[column sep=1.6em]
0 \arrow[r] & (\Hom_A(M,N))_\fp \arrow[r] \arrow[d] & (\Hom_A(A^n,N))_\fp \arrow[r] \arrow[d] & (\Hom_A(A^m,N))_\fp \arrow[d] \\
0 \arrow[r] & \Hom_{A_\fp}(M_\fp,N_\fp) \arrow[r] & \Hom_{A_\fp}(A_\fp^{n},N_\fp) \arrow[r] & \Hom_{A_\fp}(A_\fp^{m},N_\fp).
\end{tikzcd}
\end{center}
最后两个竖直箭头是同构。故
\[
(\Hom_A(M,N))_\fp\cong\Hom_{A_\fp}(M_\fp,N_\fp).
\]
类似地，对每个 $f\in A$，有典范同构
\[
(\Hom_A(M,N))_f\cong\Hom_{A_f}(M_f,N_f).
\]
定义同态
\[
\delta\colon\mathcal{H}om_{\cO_{\Spec A}}(M^\sim,N^\sim)\to(\Hom_A(M,N))^\sim
\]
如下：给定 $\mathcal{H}om_{\cO_{\Spec A}}(M^\sim,N^\sim)(U)$ 中的截面 $\phi\colon M^\sim|_U\to N^\sim|_U$，对任意 $\fp\in U$，设 $\phi_\fp\colon M_\fp\to N_\fp$ 是 $\phi$ 在茎上诱导的同态。定义映射 $\delta(\phi)\colon U\to\coprod_{\fp\in\Spec A}(\Hom_A(M,N))_\fp$，使对任意 $\fp\in U$，$(\delta(\phi))(\fp)$ 是 $\phi_\fp\colon M_\fp\to N_\fp$ 在典范同构 $\Hom_{A_\fp}(M_\fp,N_\fp)\cong(\Hom_A(M,N))_\fp$ 下的像。我们断言 $\delta(\phi)$ 是 $(\Hom_A(M,N))^\sim(U)$ 中的截面。为此，用形如 $D(f)$（$f\in A$）的开子集覆盖 $U$。则 $\phi|_{D(f)}$ 是 $\mathcal{H}om_{\cO_{\Spec A}}(M^\sim,N^\sim)(D(f))$ 中的截面。但有
\[
\mathcal{H}om_{\cO_{\Spec A}}(M^\sim,N^\sim)(D(f))
=\Hom_{\cO_{\Spec A}|_{D(f)}}(M^\sim|_{D(f)},N^\sim|_{D(f)})
\cong\Hom_{\cO_{\Spec A_f}}(M_f^\sim,N_f^\sim),
\]
\[
\Hom_{\cO_{\Spec A_f}}(M_f^\sim,N_f^\sim)\cong\Hom_{A_f}(M_f,N_f)\cong(\Hom_A(M,N))_f.
\]
故 $\phi|_{D(f)}$ 对应 $(\Hom_A(M,N))_f$ 中的一个元素。因此 $\delta(\phi)$ 是 $(\Hom_A(M,N))^\sim(U)$ 的截面。这样得到层的态射
\[
\delta\colon\mathcal{H}om_{\cO_{\Spec A}}(M^\sim,N^\sim)\to(\Hom_A(M,N))^\sim.
\]
可以验证它在 $\fp$ 处的茎上诱导同构 $\Hom_{A_\fp}(M_\fp,N_\fp)\cong(\Hom_A(M,N))_\fp$。故 $\delta$ 是同构。
\end{proof}

\begin{proposition}\label{it:1.4.3}
设 $\phi\colon A\to B$ 是环同态，$f\colon\Spec B\to\Spec A$ 是相应的态射。
\begin{enumerate}[label=(\roman*),leftmargin=2.8em]
  \item 对每个 $B$-模 $N$，有 $f_*N^\sim\cong({}_A N)^\sim$，其中右边把 $N$ 视为 $A$-模。
  \item 对每个 $A$-模 $M$，有 $f^*M^\sim\cong(B\otimes_A M)^\sim$。
\end{enumerate}
\end{proposition}

\begin{proof}
(i) 定义
\[
\alpha\colon({}_A N)^\sim\to f_*N^\sim
\]
如下：对 $({}_A N)^\sim$ 的每个截面 $s\colon U\to\coprod_{\fp\in\Spec A}({}_A N)_\fp$，定义 $(f_*N^\sim)(U)=N^\sim(f^{-1}(U))$ 的截面 $\alpha(s)\colon f^{-1}(U)\to\coprod_{\fq\in\Spec B}N_\fq$，使对每个 $\fq\in f^{-1}(U)$，$(\alpha(s))(\fq)$ 是 $s(\phi^{-1}(\fq))$ 在典范同态
\[
({}_A N)_{\phi^{-1}(\fq)}\to N_\fq,\qquad \frac{n}{s}\mapsto\frac{n}{\phi(s)}
\]
下的像。为证 $\alpha$ 是同构，只需证明：对每个 $a\in A$，同态 $({}_A N)^\sim(D(a))\xrightarrow{\alpha}(f_*N^\sim)(D(a))$ 是同构。有
\[
({}_A N)^\sim(D(a))\cong({}_A N)_a,\qquad
(f_*N^\sim)(D(a))\cong N^\sim(f^{-1}(D(a)))=N^\sim(D(\phi(a)))\cong N_{\phi(a)}.
\]
我们的断言来自以下事实：
\[
({}_A N)_a\to N_{\phi(a)},\qquad \frac{n}{a^k}\mapsto\frac{n}{\phi(a)^k}
\]
是同构。

(ii) $A$-模同态
\[
M\to{}_A(B\otimes_A M),\qquad m\mapsto1\otimes m
\]
诱导态射
\[
M^\sim\to({}_A(B\otimes_A M))^\sim\cong f_*(B\otimes_A M)^\sim.
\]
因 $f^*$ 是 $f_*$ 的左伴随，得到态射
\[
f^*M^\sim\to(B\otimes_A M)^\sim.
\]
为证它是同构，只需证明它在茎上诱导同构。对每个 $\fq\in\Spec B$，有
\[
(f^*M^\sim)_\fq\cong\cO_{\Spec B,\fq}\otimes_{\cO_{\Spec A,f(\fq)}}M^\sim_{f(\fq)}
\cong B_\fq\otimes_{A_{\phi^{-1}(\fq)}}M_{\phi^{-1}(\fq)},
\]
\[
(B\otimes_A M)^\sim_\fq\cong(B\otimes_A M)_\fq.
\]
可以验证态射 $f^*M^\sim\to(B\otimes_A M)^\sim$ 在茎上诱导典范同构
\[
B_\fq\otimes_{A_{\phi^{-1}(\fq)}}M_{\phi^{-1}(\fq)}\cong(B\otimes_A M)_\fq.
\]
\end{proof}

设 $(X,\cO_X)$ 是概形。$\cO_X$-模层 $\cF$ 称为\emph{拟凝聚}的\idx{niningju}{拟凝聚}，若 $X$ 可被仿射开子概形 $U_i=\Spec A_i$（$i\in I$）覆盖，使 $\cF|_{U_i}\cong M_i^\sim$，其中 $M_i$ 是某个 $A_i$-模。设 $X$ 是诺特概形。称 $\cF$ 为\emph{凝聚}的\idx{ningju}{凝聚}，若上述 $M_i$ 可取为有限生成。留给读者验证：对每个开子概形 $U$，若 $\cF$ 拟凝聚则 $\cF|_U$ 拟凝聚；当 $X$ 是诺特概形时，若 $\cF$ 凝聚则 $\cF|_U$ 凝聚。
\begin{tnote}
「凝聚」的定义要求 $X$ 是诺特概形，是因为要用到「有限生成模的局部化仍有限生成」以及仿射开集上的模可由有限个元素生成这类有限性。去掉诺特假设后，拟凝聚层可以局部由无限生成的模给出，凝聚性便失去意义。注意原书的凝聚定义只要求每个 $M_i$ 有限生成，并不要求核也有限生成——在诺特概形上这两件事等价（因为有限生成模的子模仍有限生成）。[依据：原书 p.66 的定义与 p.67 命题 1.4.5]
\end{tnote}

\begin{lemma}\label{it:1.4.4}
设 $X=\Spec A$ 是仿射概形，$f\in A$，$\cF$ 是拟凝聚 $\cO_X$-模。
\begin{enumerate}[label=(\roman*),leftmargin=2.8em]
  \item 对每个满足 $s|_{D(f)}=0$ 的截面 $s\in\cF(X)$，存在自然数 $n$，使 $f^ns=0$。
  \item 对每个截面 $t\in\cF(D(f))$，存在自然数 $n$，使 $f^nt$ 可延拓为 $\cF(X)$ 中的截面。
  \item 设 $M=\cF(X)$。则 $M_f\cong\cF(D(f))$。
\end{enumerate}
\end{lemma}

\begin{proof}
用有限多个仿射开子概形 $U_i=\Spec A_i$（$i\in I$）覆盖 $\Spec A$，使 $\cF|_{U_i}=M_i^\sim$，其中 $M_i$ 是某个 $A_i$-模。对每个 $i$，设 $f_i\in A_i$ 是 $f$ 在限制映射
\[
A\cong\cO(\Spec A)\to\cO(\Spec A_i)\cong A_i
\]
下的像。则 $D(f)\cap U_i$ 与 $U_i=\Spec A_i$ 的开子集 $D(f_i)$ 重合。故
\[
\cF(D(f)\cap U_i)\cong M_i^\sim(D(f_i))\cong(M_i)_{f_i}.
\]
设 $s\in\cF(X)$ 满足 $s|_{D(f)}=0$。则对每个 $i$，$s|_{U_i}$ 落在限制映射 $\cF(U_i)\to\cF(D(f)\cap U_i)$ 的核中。这个限制映射可等同于典范同态 $M_i\to(M_i)_{f_i}$。因 $M_i\to(M_i)_{f_i}$ 的核中每个元素都被 $f_i$ 的某个幂零化，故对某个自然数 $n_i$ 有 $f_i^{n_i}s|_{U_i}=0$。因覆盖 $\{U_i\}_{i\in I}$ 是有限覆盖，存在自然数 $n$，使 $f^ns|_{U_i}=0$ 对每个 $i$ 成立。故 $f^ns=0$。这证明了 (i)。

设 $t\in\cF(D(f))$。则 $t|_{D(f)\cap U_i}\in\cF(D(f)\cap U_i)$。因限制映射 $\cF(U_i)\to\cF(D(f)\cap U_i)$ 可等同于典范同态 $M_i\to(M_i)_{f_i}$，存在自然数 $n$，使 $f_i^nt|_{D(f)\cap U_i}$ 对每个 $i$ 都可延拓为 $\cF(U_i)$ 中的截面 $t_i$。注意对每一对 $i,j$，有
\[
t_i|_{D(f)\cap U_i\cap U_j}=t_j|_{D(f)\cap U_i\cap U_j}.
\]
此外，由 \ref{it:1.3.25} (i)，$U_i\cap U_j$ 是仿射的。对拟凝聚层 $\cF|_{U_i\cap U_j}$ 应用 (i)，得到自然数 $m$，使对每一对 $i,j$ 有
\[
f_i^mt_i|_{U_i\cap U_j}=f_j^mt_j|_{U_i\cap U_j}.
\]
于是可以把 $f_i^mt_i$ 粘成 $\cF(X)$ 中的一个截面。它是 $f^{m+n}t$ 的延拓。这证明了 (ii)。
\end{proof}

有一个同态 $M_f\to\cF(D(f))$，由 $\frac{s}{f^k}\mapsto\frac{1}{f^k}\cdot s|_{D(f)}$（$s\in M=\cF(X)$）给出。由 (i) 它是单射，由 (ii) 它是满射。故它是同构。

\begin{proposition}\label{it:1.4.5}
设 $X$ 是概形，$\cF$ 是 $X$ 上的拟凝聚层。则对 $X$ 的任意仿射开子概形 $U=\Spec A$，存在 $A$-模 $M$，使 $\cF|_U\cong M^\sim$。若 $X$ 是诺特概形且 $\cF$ 凝聚，则 $M$ 有限生成。
\end{proposition}

\begin{proof}
把 $X$ 换为 $U$、$\cF$ 换为 $\cF|_U$，可设 $X=\Spec A$。令 $M=\cF(X)$。如下定义同构 $\phi\colon M^\sim\to\cF$：对任意 $f\in A$，$\phi(D(f))$ 由
\[
\phi(D(f))\colon M^\sim(D(f))\cong M_f\to\cF(D(f)),\qquad \frac{s}{f^n}\mapsto\frac{1}{f^n}s|_{D(f)}
\]
给出，其中 $s\in M\cong\cF(X)$。由 \ref{it:1.4.4} (iii)，$\phi(D(f))$ 是同构。这就对每个形如 $V=D(f)$ 的开子集定义了 $\phi(V)\colon M^\sim(V)\to\cF(V)$。一般地，用开子集 $D(f_i)$（$i\in I$）覆盖 $V$。有正合序列
\[
0\to\cF(V)\to\prod_{i\in I}\cF(D(f_i))\to\prod_{i,j\in I}\cF(D(f_i)\cap D(f_j)),
\]
其中第二个箭头由 $s\mapsto(s|_{D(f_i)})$ 定义，第三个箭头由 $(s_i)\mapsto(s_j|_{D(f_i)\cap D(f_j)}-s_i|_{D(f_i)\cap D(f_j)})$ 定义。对 $M^\sim$ 有类似的正合序列。我们定义 $\phi(V)\colon M^\sim(V)\to\cF(V)$，使下图交换：
\begin{center}
\begin{tikzcd}[column sep=1.1em,row sep=2.2em]
0 \arrow[r] & M^\sim(V) \arrow[r] \arrow[d,"\phi(V)"'] &
\prod_{i\in I}M^\sim(D(f_i)) \arrow[r] \arrow[d,"\prod_i\phi(D(f_i))"'] &
\prod_{i,j\in I}M^\sim(D(f_i)\cap D(f_j)) \arrow[d,"\prod_{i,j}\phi(D(f_i)\cap D(f_j))"] \\
0 \arrow[r] & \cF(V) \arrow[r] & \prod_{i\in I}\cF(D(f_i)) \arrow[r] &
\prod_{i,j\in I}\cF(D(f_i)\cap D(f_j)).
\end{tikzcd}
\end{center}
可以验证 $\phi(V)$ 与 $V$ 的开覆盖的选取无关。事实上，使上图交换的 $\phi(V)$ 只有一个。若 $\{D(f_i)\}_{i\in I}$ 与 $\{D(f'_{i'})\}_{i'\in I'}$ 是 $V$ 的两个覆盖，则由这两个覆盖定义的 $\phi(V)$ 都与由覆盖 $D(f_if'_{i'})$（$i\in I$，$i'\in I'$）定义的 $\phi(V)$ 一致。

设 $A$ 是诺特环且 $\cF$ 凝聚。则可用有限多个仿射开子概形 $U_i=\Spec A_i$（$i\in I$）覆盖 $\Spec A$，使 $\cF|_{U_i}\cong M_i^\sim$，其中 $M_i$ 是有限生成 $A_i$-模。用有限多个开子集 $D(f_{ij})$（$j\in J_i$）覆盖每个 $U_i$，其中 $f_{ij}\in A$。记 $f_{ij}$ 在限制映射
\[
A\cong\cO_X(\Spec A)\to\cO_X(U_i)\cong A_i
\]
下的像为 $f'_{ij}$。则
\[
\cF(D(f_{ij}))\cong M_{f_{ij}},\qquad
\cF(D(f_{ij}))\cong\cF|_{U_i}(D(f'_{ij}))\cong M_i^\sim(D(f'_{ij}))\cong(M_i)_{f'_{ij}}.
\]
故 $M_{f_{ij}}\cong(M_i)_{f'_{ij}}$。类似地有 $A_{f_{ij}}\cong(A_i)_{f'_{ij}}$。因每个 $M_i$ 是有限生成 $A_i$-模，每个 $(M_i)_{f'_{ij}}$ 是有限生成 $(A_i)_{f'_{ij}}$-模。故每个 $M_{f_{ij}}$ 是有限生成 $A_{f_{ij}}$-模。此外，$D(f_{ij})$（$i\in I$，$j\in J_i$）构成 $\Spec A$ 的开覆盖。这些事实说明 $M$ 是有限生成 $A$-模。（对照 \ref{it:1.3.11} 的证明。）
\end{proof}

\begin{corollary}\label{it:1.4.6}
设 $X$ 是概形，$\cF$ 与 $\cG$ 是拟凝聚 $\cO_X$-模。
\begin{enumerate}[label=(\roman*),leftmargin=2.8em]
  \item $\cF\otimes_{\cO_X}\cG$ 拟凝聚。若 $\cF$ 有限展示，则 $\mathcal{H}om_{\cO_X}(\cF,\cG)$ 拟凝聚。
  \item 设 $\phi\colon\cF\to\cG$ 是态射。则 $\ker\phi$、$\coker\phi$ 与 $\Image\phi$ 都拟凝聚。
\end{enumerate}
若 $X$ 是诺特概形，则对凝聚层有同样的结论。
\end{corollary}

\begin{proof}
可设 $X=\Spec A$ 是仿射的。则 $\cF\cong M^\sim$ 且 $\cG\cong N^\sim$，其中 $M$ 与 $N$ 是某个 $A$-模。由 \ref{it:1.4.2} (iv)，有
\[
\cF\otimes_{\cO_X}\cG\cong(M\otimes_A N)^\sim.
\]
故 $\cF\otimes_{\cO_X}\cG$ 拟凝聚。若 $\cF$ 有限展示，则通过缩小 $X$ 并用 \ref{it:1.4.2} (iii)--(iv)，可设 $M$ 是有限展示的 $A$-模。由 \ref{it:1.4.2} (iv)，有
\[
\mathcal{H}om_{\cO_X}(\cF,\cG)\cong(\Hom_A(M,N))^\sim.
\]
故 $\mathcal{H}om_{\cO_X}(\cF,\cG)$ 拟凝聚。设 $\phi\colon\cF\to\cG$ 是态射。由 \ref{it:1.4.2} (iv)，它由某个同态 $\psi\colon M\to N$ 诱导。由 \ref{it:1.4.2} (iii)，$\ker\phi$、$\coker\phi$ 与 $\Image\phi$ 分别同构于 $(\ker\psi)^\sim$、$(\coker\psi)^\sim$ 与 $(\Image\psi)^\sim$。故它们拟凝聚。当 $X$ 是诺特概形且 $\cF$ 与 $\cG$ 凝聚时，上述所有模都有限生成，从而定义凝聚层。
\end{proof}

$A$-模 $P$ 称为\emph{投射}的\idx{toushe}{投射模}，若对 $A$-模的每个满同态 $M\to N$，每个同态 $P\to N$ 都可提升为同态 $P\to M$。设
\[
0\to M'\to M\xrightarrow{p}M''\to0
\]
是 $A$-模的正合序列。若 $M''$ 投射，则该正合序列分裂。事实上，可把同态 $\id_{M''}$ 提升为同态 $i\colon M''\to M$。有 $pi=\id_{M''}$。我们的断言便由 \ref{it:1.1.3} 得出。对任意投射模 $P$，可找到满同态 $\phi\colon F\to P$，其中 $F$ 是自由 $A$-模。正合序列
\[
0\to\ker\phi\to F\to P\to0
\]
分裂。故 $P$ 是自由模的直和项。反之，可以证明自由模的任意直和项都是投射的。

\begin{proposition}\label{it:1.4.7}
设 $M$ 是有限展示的 $A$-模。则 $M$ 是投射 $A$-模，当且仅当 $M^\sim$ 是局部自由 $\cO_{\Spec A}$-模。
\end{proposition}

\begin{proof}
设 $M^\sim$ 局部自由。给定 $A$-模的任意正合序列
\[
N\to N'\to0,
\]
有 $\cO_{\Spec A}$-模的正合序列
\[
N^\sim\to N'^\sim\to0.
\]
因 $M^\sim$ 局部自由且其秩必有限，序列
\[
\mathcal{H}om_{\cO_{\Spec A}}(M^\sim,N^\sim)\to\mathcal{H}om_{\cO_{\Spec A}}(M^\sim,N'^\sim)\to0
\]
正合。（为此可局部化，从而设 $M^\sim$ 是有限秩自由模。此时断言显然。）故由 \ref{it:1.4.2} (iv)，序列
\[
(\Hom_A(M,N))^\sim\to(\Hom_A(M,N'))^\sim\to0
\]
正合。由 \ref{it:1.4.2} (iii)，序列
\[
\Hom_A(M,N)\to\Hom_A(M,N')\to0
\]
正合。故 $M$ 投射。

反之，设 $M$ 投射，即 $M$ 是自由 $A$-模的直和项。则对每个 $\fp\in\Spec A$，$M_\fp$ 也是自由 $A_\fp$-模的直和项，从而投射。由下面的 \ref{it:1.4.8}，$M_\fp$ 自由。故有同构 $M_\fp\cong A_\fp^n$，其中 $n$ 是某个自然数。由 \ref{it:1.4.1} (ii)，存在 $\fp$ 的邻域 $U$，使 $M^\sim|_U\cong\cO_{\Spec A}|_U$。故 $M^\sim$ 局部自由。
\end{proof}

\begin{lemma}\label{it:1.4.8}
设 $A$ 是局部环，$M$ 是有限型 $A$-模。则 $M$ 投射，当且仅当 $M$ 自由。
\end{lemma}

\begin{proof}
设 $M$ 是有限型投射 $A$-模，我们来证明 $M$ 自由。设 $\fm$ 是 $A$ 的极大理想。选取 $x_1,\dots,x_n$，使它们在 $M/\fm M$ 中的像构成 $A/\fm$-向量空间 $M/\fm M$ 的一组基。考虑同态
\[
A^n\to M,\qquad e_i\mapsto x_i\quad(i=1,\dots,n),
\]
其中对每个 $i$，$e_i$ 是 $A^n$ 中第 $i$ 个分量为 $1$、其余分量为 $0$ 的元素。由 Nakayama 引理，$A^n\to M$ 是满射。设 $R$ 是其核。因 $M$ 投射，序列
\[
0\to R\to A^n\to M\to0
\]
分裂正合。这说明
\[
0\to R/\fm R\to(A/\fm A)^n\to M/\fm M\to0
\]
分裂正合。由 $A^n\to M$ 的选取，同态 $(A/\fm)^n\to M/\fm M$ 是同构。故 $R/\fm R=0$。因 $R$ 是 $A^n$ 的直和项，它有限型。由 Nakayama 引理，$R=0$。故 $A^n\to M$ 是同构，从而 $M$ 自由。
\end{proof}
\begin{tnote}
Nakayama 引理的标准形式：设 $A$ 是局部环，$\fm$ 是其极大理想，$M$ 是有限生成 $A$-模；若 $\fm M=M$，则 $M=0$。上面的论证两次用到它：先由 $M/\fm M$ 有限维得 $A^n\to M$ 满，再由 $R/\fm R=0$ 与 $R$ 有限型得 $R=0$。[来源：Atiyah--Macdonald, \emph{Introduction to Commutative Algebra}, Proposition 2.8]
\end{tnote}

\begin{proposition}\label{it:1.4.9}
设 $f\colon X\to Y$ 是概形态射。
\begin{enumerate}[label=(\roman*),leftmargin=2.8em]
  \item 若 $\cG$ 是拟凝聚 $\cO_Y$-模，则 $f^*\cG$ 是拟凝聚 $\cO_X$-模。
  \item 若 $X$ 与 $Y$ 是诺特概形且 $\cG$ 凝聚，则 $f^*\cG$ 凝聚。
  \item 若 $f$ 拟紧且拟分离，$\cF$ 是拟凝聚 $\cO_X$-模，则 $f_*\cF$ 是拟凝聚 $\cO_Y$-模。
\end{enumerate}
\end{proposition}

\begin{proof}
问题关于 $Y$ 是局部的。可设 $Y=\Spec A$ 是仿射的。用仿射开子概形 $U_i=\Spec B_i$（$i\in I$）覆盖 $X$。对每个 $i$，记 $f_i$ 为 $f$ 诱导的态射 $U_i\to Y$。

(i) 若 $\cG$ 是拟凝聚 $\cO_Y$-模，则 $\cG\cong M^\sim$，其中 $M$ 是某个 $A$-模。对每个 $i$，由 \ref{it:1.4.3} (ii) 有
\[
f^*\cG|_{U_i}\cong f_i^*\cG\cong(B_i\otimes_A M)^\sim.
\]
故 $f^*\cG$ 拟凝聚。

(ii) 若 $X$ 与 $Y$ 是诺特概形且 $\cG$ 凝聚，则 $M$ 是有限生成 $A$-模，从而每个 $B_i\otimes_A M$ 是有限生成 $B_i$-模。故 $f^*\cG$ 凝聚。

(iii) 若 $f$ 拟紧且拟分离，则可设 $\{U_i\}_{i\in I}$ 是 $X$ 的有限覆盖，且由 \ref{it:1.3.25} (ii)，每个 $U_i\cap U_j$ 可被有限多个仿射开子概形 $U_{ijk}$（$k\in K_{ij}$）覆盖。记 $f_{ijk}$ 为 $f$ 诱导的态射 $U_{ijk}\to Y$。对 $Y$ 的任意开子集 $V$，有正合序列
\[
0\to\cF(f^{-1}(V))\to\prod_{i\in I}\cF(f^{-1}(V)\cap U_i)\to\prod_{i,j\in I,\,k\in K_{ij}}\cF(f^{-1}(V)\cap U_{ijk}),
\]
其中第二个箭头由 $s\mapsto(s|_{f^{-1}(V)\cap U_i})$ 定义，第三个箭头由 $(s_i)\mapsto(s_j|_{f^{-1}(V)\cap U_{ijk}}-s_i|_{f^{-1}(V)\cap U_{ijk}})$ 定义。故有正合序列
\[
0\to f_*\cF\to\prod_{i\in I}f_{i*}(\cF|_{U_i})\to\prod_{i,j\in I,\,k\in K_{ij}}f_{ijk*}(\cF|_{U_{ijk}}).
\]
由 \ref{it:1.4.2} (v) 与 \ref{it:1.4.3} (i)，$\prod_{i\in I}f_{i*}(\cF|_{U_i})$ 与 $\prod_{i,j\in I,\,k\in K_{ij}}f_{ijk*}(\cF|_{U_{ijk}})$ 拟凝聚。故 $f_*\cF$ 拟凝聚。
\end{proof}

\begin{proposition}\label{it:1.4.10}
设 $f\colon X\to S$ 是仿射态射，$\cF$ 与 $\cG$ 是拟凝聚 $\cO_X$-模。
\begin{enumerate}[label=(\roman*),leftmargin=2.8em]
  \item 有双射
  \[
  \Hom_{\cO_X}(\cF,\cG)\cong\Hom_{f_*\cO_X}(f_*\cF,f_*\cG),
  \]
  其中右边是 $f_*\cO_X$-模的态射集合。
  \item 设 $\phi\colon\cF\to\cG$ 是态射。则 $\phi$ 是同构，当且仅当 $\phi$ 诱导的 $f_*\cO_X$-模态射 $f_*\phi\colon f_*\cF\to f_*\cG$ 是同构。
\end{enumerate}
\end{proposition}

\begin{proof}
(i) 任意 $\cO_X$-模态射 $\phi\colon\cF\to\cG$ 诱导 $f_*\cO_X$-模态射 $f_*\phi\colon f_*\cF\to f_*\cG$。定义
\[
F\colon\Hom_{\cO_X}(\cF,\cG)\to\Hom_{f_*\cO_X}(f_*\cF,f_*\cG)
\]
为 $F(\phi)=f_*\phi$。我们证明它是双射。先考虑 $S$ 仿射的情形。则 $X$ 也仿射。由 \ref{it:1.4.3} (i)，$f_*\cF$ 与 $f_*\cG$ 拟凝聚。由 \ref{it:1.4.2} (iv)，有
\[
\Hom_{\cO_X}(\cF,\cG)\cong\Hom_{\cO_X(X)}(\cF(X),\cG(X)),
\]
\[
\Hom_{\cO_S}(f_*\cF,f_*\cG)\cong\Hom_{\cO_S(S)}((f_*\cF)(S),(f_*\cG)(S)).
\]
群 $\Hom_{f_*\cO_X}(f_*\cF,f_*\cG)$ 由 $\Hom_{\cO_S}(f_*\cF,f_*\cG)$ 中使下图交换的 $\cO_S$-模同态 $\psi\colon f_*\cF\to f_*\cG$ 组成：
\begin{center}
\begin{tikzcd}[column sep=3em,row sep=2.4em]
f_*\cO_X\otimes_{\cO_S}f_*\cF \arrow[r,"\id\otimes\psi"] \arrow[d] & f_*\cO_X\otimes_{\cO_S}f_*\cG \arrow[d] \\
f_*\cF \arrow[r,"\psi"'] & f_*\cG,
\end{tikzcd}
\end{center}
其中竖直箭头来自 $f_*\cF$ 与 $f_*\cG$ 上的 $(f_*\cO_X)$-模结构。因上图中的层都是拟凝聚 $\cO_S$-模，上图交换当且仅当以下图交换：
\begin{center}
\begin{tikzcd}[column sep=3em,row sep=2.4em]
(f_*\cO_X)(S)\otimes_{\cO_S(S)}(f_*\cF)(S) \arrow[r,"\id\otimes\psi"] \arrow[d] &
(f_*\cO_X)(S)\otimes_{\cO_S(S)}(f_*\cG)(S) \arrow[d] \\
(f_*\cF)(S) \arrow[r,"\psi"'] & (f_*\cG)(S).
\end{tikzcd}
\end{center}
于是有
\[
\Hom_{f_*\cO_X}(f_*\cF,f_*\cG)\cong\Hom_{(f_*\cO_X)(S)}((f_*\cF)(S),(f_*\cG)(S))\cong\Hom_{\cO_X(X)}(\cF(X),\cG(X)).
\]
由此可见
\[
\Hom_{\cO_X}(\cF,\cG)\cong\Hom_{f_*\cO_X}(f_*\cF,f_*\cG).
\]
这证明了仿射情形下的断言。一般地，用仿射开子概形 $V_i$（$i\in I$）覆盖 $S$，并用仿射开子概形 $V_{ijk}$（$k\in K_{ij}$）覆盖每个 $V_i\cap V_j$。有交换图
\begin{center}\small
\begin{tikzcd}[column sep=1.8em,row sep=2.6em]
\prod_{i\in I}\Hom_{\cO_X|_{f^{-1}(V_i)}}\bigl(\cF|_{f^{-1}(V_i)},\cG|_{f^{-1}(V_i)}\bigr)
\arrow[r] \arrow[d] &
\prod_{i,j\in I,\,k\in K_{ij}}\Hom_{\cO_X|_{f^{-1}(V_{ijk})}}\bigl(\cF|_{f^{-1}(V_{ijk})},\cG|_{f^{-1}(V_{ijk})}\bigr) \arrow[d] \\
\prod_{i\in I}\Hom_{(f_*\cO_X)|_{V_i}}\bigl((f_*\cF)|_{V_i},(f_*\cG)|_{V_i}\bigr)
\arrow[r] &
\prod_{i,j\in I,\,k\in K_{ij}}\Hom_{(f_*\cO_X)|_{V_{ijk}}}\bigl((f_*\cF)|_{V_{ijk}},(f_*\cG)|_{V_{ijk}}\bigr),
\end{tikzcd}
\end{center}
其中水平箭头分别由
\[
(\phi_i)\mapsto\bigl(\phi_j|_{f^{-1}(V_{ijk})}-\phi_i|_{f^{-1}(V_{ijk})}\bigr),\qquad
(\phi_i)\mapsto\bigl(\phi_j|_{V_{ijk}}-\phi_i|_{V_{ijk}}\bigr)
\]
定义。由上面已证的仿射情形，两个竖直箭头都是同构。因 $\Hom_{\cO_X}(\cF,\cG)$ 与 $\Hom_{f_*\cO_X}(f_*\cF,f_*\cG)$ 可等同于水平箭头的核，有
\[
\Hom_{\cO_X}(\cF,\cG)\cong\Hom_{f_*\cO_X}(f_*\cF,f_*\cG).
\]

(ii) 我们证明「当」的部分。设 $f_*\phi\colon f_*\cF\to f_*\cG$ 是同构。由 (i)，存在态射 $\psi\colon\cG\to\cF$，使 $f_*\psi$ 是 $f_*\phi$ 的逆。有
\[
f_*(\psi\phi)=\id_{f_*\cF}=f_*(\id_{\cF}).
\]
故由 (i) 得 $\psi\phi=\id_{\cF}$。类似地 $\phi\psi=\id_{\cG}$。故 $\phi$ 是同构。
\end{proof}

\begin{proposition}\label{it:1.4.11}
设 $f\colon X\to S$ 与 $g\colon Y\to S$ 是概形态射。设 $g$ 是仿射的。则有双射
\[
\Hom_S(X,Y)\cong\Hom_{\cO_S}(g_*\cO_Y,f_*\cO_X),
\]
其中左边是 $S$-态射 $X\to Y$ 的集合，右边是 $\cO_S$-代数层从 $g_*\cO_Y$ 到 $f_*\cO_X$ 的态射集合。
\end{proposition}

\begin{proof}
定义
\[
F\colon\Hom_S(X,Y)\to\Hom_{\cO_S}(g_*\cO_Y,f_*\cO_X),
\]
使对每个 $S$-态射 $h\colon X\to Y$，$F(h)$ 为复合
\[
g_*\cO_Y\xrightarrow{g_*(h^\sharp)}g_*h_*\cO_X\cong f_*\cO_X.
\]
固定 $S$ 的仿射开子概形覆盖 $\{V_i\}_{i\in I}$。则每个 $g^{-1}(V_i)$ 是仿射的。故由 \ref{it:1.2.8}，有
\[
\Hom_{V_i}\bigl(f^{-1}(V_i),g^{-1}(V_i)\bigr)\cong\Hom_{\cO_S(V_i)}\bigl(\cO_Y(g^{-1}(V_i)),\cO_X(f^{-1}(V_i))\bigr).
\]
定义映射
\[
G\colon\Hom_{\cO_S}(g_*\cO_Y,f_*\cO_X)\to\Hom_S(X,Y)
\]
如下：给定 $\cO_S$-代数态射 $\phi\colon g_*\cO_Y\to f_*\cO_X$，对每个 $i$，定义 $h_i\colon f^{-1}(V_i)\to g^{-1}(V_i)$ 为与 $\cO_S(V_i)$-代数同态
\[
\cO_Y(g^{-1}(V_i))=(g_*\cO_Y)(V_i)\xrightarrow{\phi}(f_*\cO_X)(V_i)=\cO_X(f^{-1}(V_i))
\]
相应的 $V_i$-态射。可以验证 $h_i|_{f^{-1}(V_i)\cap f^{-1}(V_j)}=h_j|_{f^{-1}(V_i)\cap f^{-1}(V_j)}$。故可把 $h_i$ 粘成 $S$-态射 $h\colon X\to Y$。定义 $G(\phi)=h$。可以验证 $F$ 与 $G$ 互为逆。
\end{proof}

设 $f\colon X\to S$ 是仿射态射。由 \ref{it:1.4.3} (i)，$\cO_S$-代数层 $f_*\cO_X$ 是拟凝聚 $\cO_S$-模。反之，有如下命题。

\begin{proposition}\label{it:1.4.12}
设 $S$ 是概形，$B$ 是 $\cO_S$-代数层，且它作为 $\cO_S$-模拟凝聚。（称这样的 $B$ 为拟凝聚 $\cO_S$-代数。）则存在仿射态射 $f\colon X\to S$，在唯一同构意义下唯一，使 $f_*\cO_X\cong B$。此外，对每个作为 $\cO_S$-模拟凝聚的 $B$-模层 $\cF$，存在拟凝聚 $\cO_X$-模 $\cF'$，在唯一同构意义下唯一，使 $f_*\cF'\cong\cF$ 作为 $B$-模层。
\end{proposition}

我们常把仿射 $S$-概形 $X$ 记作 $\Spec B$。

\begin{proof}
设 $f_1\colon X_1\to S$ 与 $f_2\colon X_2\to S$ 是两个仿射态射，使 $f_{i*}\cO_{X_i}\cong B$。由 \ref{it:1.4.11}，存在唯一态射 $g\colon X_1\to X_2$，使 $f_2g=f_1$ 且下图交换：
\begin{center}
\begin{tikzcd}[column sep=3.2em,row sep=2.4em]
f_{2*}\cO_{X_2} \arrow[r,"f_{2*}(g^\sharp)"] \arrow[dr,"\cong"'] & f_{1*}\cO_{X_1} \arrow[d,"\cong"] \\
& B.
\end{tikzcd}
\end{center}
类似地，存在唯一态射 $h\colon X_2\to X_1$，使 $f_1h=f_2$ 且下图交换：
\begin{center}
\begin{tikzcd}[column sep=3.2em,row sep=2.4em]
f_{1*}\cO_{X_1} \arrow[r,"f_{1*}(h^\sharp)"] \arrow[dr,"\cong"'] & f_{2*}\cO_{X_2} \arrow[d,"\cong"] \\
& B.
\end{tikzcd}
\end{center}
再用 \ref{it:1.4.11} 可以验证 $hg=\id_{X_1}$ 且 $gh=\id_{X_2}$。这证明了满足 $f_*\cO_X\cong B$ 的仿射态射 $f\colon X\to S$ 的唯一性。类似地，用 \ref{it:1.4.10} (i) 可证明满足 $f_*\cF'\cong\cF$ 的拟凝聚 $\cO_X$-模 $\cF'$ 的唯一性。

用仿射开子概形 $V_i=\Spec A_i$（$i\in I$）覆盖 $S$。有 $B|_{V_i}\cong B_i^\sim$，其中 $B_i$ 是某个 $A_i$-模。由 \ref{it:1.4.2} (iv)，给出 $\cO_S$-代数 $B$ 的乘法的态射 $B\otimes_{\cO_S}B\to B$ 决定同态 $B_i\otimes_{A_i}B_i\to B_i$。这些同态定义了 $B_i$ 上的乘法，它们交换且结合。故 $B_i$ 是 $A_i$-代数。类似地有 $\cF|_{V_i}\cong M_i^\sim$，其中 $M_i$ 是某个 $B_i$-模。对每个 $i$，设 $X_i=\Spec B_i$，设 $\cF_i'$ 为由 $B_i$-模 $M_i$ 定义的 $\cO_{X_i}$-模，并设 $f_i\colon X_i\to V_i$ 为由 $B_i$ 上的 $A_i$-代数结构定义的态射。有
\[
f_{i*}\cO_{X_i}\cong B|_{V_i},\qquad f_{i*}\cF_i'\cong\cF|_{V_i}.
\]
故对每一对 $i,j$，有
\[
(f_{i*}\cO_{X_i})|_{V_i\cap V_j}\cong(f_{j*}\cO_{X_j})|_{V_i\cap V_j},\qquad
(f_{i*}\cF_i')|_{V_i\cap V_j}\cong(f_{j*}\cF_j')|_{V_i\cap V_j}.
\]
由 \ref{it:1.4.11}，有
\[
f_i^{-1}(V_i\cap V_j)\cong f_j^{-1}(V_i\cap V_j)
\]
作为 $V_i\cap V_j$ 上的概形，且由 \ref{it:1.4.10} 有
\[
\cF_i'|_{f_i^{-1}(V_i\cap V_j)}\cong\cF_j'|_{f_j^{-1}(V_i\cap V_j)}.
\]
留给读者证明这些同构给出 $X_i$ 与 $\cF_i'$ 的粘合数据。把它们粘起来，便得到具有所需性质的 $X$ 与 $\cF'$。
\end{proof}

设 $(X,\cO_X)$ 是环层空间。$\cO_X$ 的\emph{理想层}\idx{lixiangceng}{理想层} $\cI$ 是 $\cO_X$ 的 $\cO_X$-子模。对每个 $\cO_X$-模 $\cF$，定义 $\cI\cF$ 为态射
\[
\cI\otimes_{\cO_X}\cF\to\cO_X\otimes_{\cO_X}\cF\cong\cF
\]
的像。

设 $(f,f^\sharp)\colon(X,\cO_X)\to(Y,\cO_Y)$ 是环层空间态射，$\cJ$ 是 $\cO_Y$ 的理想层。定义 $\cJ\cO_X$ 为态射
\[
f^*\cJ\to f^*\cO_Y\cong\cO_X
\]
的像。它是 $\cO_X$ 的理想层。更一般地，对每个 $\cO_X$-模 $\cF$，定义 $\cJ\cF$ 为态射
\[
f^*\cJ\otimes_{\cO_X}\cF\to f^*\cO_Y\otimes_{\cO_X}\cF\cong\cF
\]
的像。有
\[
\cJ\cF=(\cJ\cO_X)\cF.
\]

\begin{proposition}\label{it:1.4.13}
设 $(f,f^\sharp)\colon(X,\cO_X)\to(Y,\cO_Y)$ 是环层空间态射，$\cJ$ 是 $\cO_Y$ 的理想层。对每个 $\cO_X$-模 $\cF$，有
\[
\cF\otimes_{\cO_X}f^*(\cO_Y/\cJ)\cong\cF/\cJ\cF.
\]
\end{proposition}

\begin{proof}
有正合序列
\[
0\to\cJ\to\cO_Y\to\cO_Y/\cJ\to0,
\]
从而有正合序列
\[
0\to f^{-1}\cJ\to f^{-1}\cO_Y\to f^{-1}(\cO_Y/\cJ)\to0.
\]
用环层 $f^{-1}\cO_Y$ 与 $\cF\otimes_{\cO_X}\cO_X$ 作张量积，得到正合序列
\[
\cF\otimes_{\cO_X}(\cO_X\otimes_{f^{-1}\cO_Y}f^{-1}\cJ)\to
\cF\otimes_{\cO_X}(\cO_X\otimes_{f^{-1}\cO_Y}f^{-1}\cO_Y)
\]
\[
\to\cF\otimes_{\cO_X}(\cO_X\otimes_{f^{-1}\cO_Y}f^{-1}(\cO_Y/\cJ))\to0,
\]
即
\[
\cF\otimes_{\cO_X}f^*\cJ\to\cF\to\cF\otimes_{\cO_X}f^*(\cO_Y/\cJ)\to0.
\]
这说明
\[
\cF/\cJ\cF\cong\cF\otimes_{\cO_X}f^*(\cO_Y/\cJ).
\]
\end{proof}
